% !TeX program = xelatex
% UTF-8; standalone XeLaTeX + XePersian document.
% Source: Google Drive / notebook / ped neurological evaluation.
% pic01.jpg = 20261016_210538757.jpg (1huZDVwNgOw-eV0oM9pNa_DU9aEHIzmUh)
% pic02.jpg = 20261016_210544565.jpg (1AgHoelvsMSnfk3ioAG3QIa1cnquF92D0)
% pic03.jpg = 20261016_210548657.jpg (1jwh03TfnCwXV-iCwL54gU5KzW_jtgmge)
% pic04.jpg = 20261016_210554543.jpg (18LH4MYPpGOxSYwebWDd2K0VzrSnFxDjR)
% pic05.jpg = 20261016_210559216.jpg (1Yw_2Us935PI923k8BlNrQQFkZF9Zl25p)
% pic06.jpg = 20261016_210603428.jpg (1Ymd-AyOcamAdKVr6zElkfBZ2YAltzt3y)
% Dual-mode guard: compiles standalone with XeLaTeX, or is \input from the
% XePersian notebook (main.tex) where its local preamble below is skipped.
\newif\ifPedNeuroEvalStandalone
\ifdefined\settextfont
  \PedNeuroEvalStandalonefalse
\else
  \PedNeuroEvalStandalonetrue
\fi
\ifPedNeuroEvalStandalone
\documentclass[11pt,a4paper]{article}
\usepackage[margin=20mm]{geometry}
\usepackage{amsmath,amssymb}
\usepackage{enumitem}
\usepackage{xcolor}
\usepackage{xepersian}
\settextfont{Amiri}
\setlatintextfont{DejaVu Serif}
\setlength{\parindent}{0pt}
\setlength{\parskip}{4pt}
\setlength{\emergencystretch}{3em}
\setlist[itemize]{topsep=3pt,itemsep=3pt}
\newcommand{\en}[1]{\lr{#1}}
\newcommand{\marktext}[1]{\colorbox{yellow!30}{\strut #1}}
\newcommand{\uncertain}[1]{\textcolor{red!65!black}{[خوانش نامطمئن: #1]}}
\newcommand{\pagehead}[1]{%
  \section*{تصویر #1}
  \begingroup\small
  \en{Subject:} ۱۷۹\quad
  \en{Year:}\quad\en{Month:}\quad\en{Date:}
  (خانه‌های تاریخ خالی‌اند). شماره‌های چاپی سطرها: ۱ تا ۹.
  \par\endgroup}
\begin{document}
\fi

% Helper macros used inside the body must exist in notebook mode too.
% (Each \providecommand must carry a full body: a body-less \providecommand
% consumes the following token as its definition and corrupts the document.)
\providecommand{\marktext}[1]{\colorbox{yellow!30}{\strut #1}}
\providecommand{\uncertain}[1]{\textcolor{red!65!black}{[خوانش نامطمئن: #1]}}
\providecommand{\pagehead}[1]{%
  \section*{تصویر #1}
  \begingroup\small
  \en{Subject:} ۱۷۹\quad
  \en{Year:}\quad\en{Month:}\quad\en{Date:}
  (خانه‌های تاریخ خالی‌اند). شماره‌های چاپی سطرها: ۱ تا ۹.
  \par\endgroup}

\begin{center}
{\Large\bfseries ۱۷۹ ـ ارزیابی‌های نورولوژیک}
\end{center}

\pagehead{۱}
\subsection*{الگوی حملات و پسرفت}
\begin{latin}
\begin{itemize}
  \item Intermittent attacks of recurrent, stereotyped episodes
  $\rightarrow$ Epilepsy, Migraine, Vascular syndromes.
  \item Loss of skills over time (Regression):
  \begin{itemize}
    \item Underlying degenerative diseases,
    \item Emerging autism spectrum disorders,
    \item Epileptic encephalopathy.
  \end{itemize}
\end{itemize}
\end{latin}

\subsection*{یافته‌های پوست و سر}
\begin{itemize}
  \item ماکول‌های \marktext{\en{Café-au-lait}}: \en{flat}، قهوه‌ای کم‌رنگ؛ \marktext{\en{NF I}}.
  \item در توبروس اسکلروزیس:
  \begin{itemize}
    \item \en{Facial angiofibromas}: مانند آکنه روی بینی و ناحیهٔ مالار.
    \item \en{Nail fibromas}.
    \item \en{\marktext{Ash-leaf spots} (hypopigmented macules)}.
    \item \en{\marktext{Shagreen patches} (flesh-colored soft plaques with prominent follicular opening)}.
  \end{itemize}
  \item فونتانل برجسته: افزایش \en{ICP}؛ گریه و وضعیت نیز می‌تواند باعث برجسته شدن بشود.
  \item کرانیوسینوستوز: زود بسته شدن سوچورها؛ در سندرم‌های ژنتیکی.
\end{itemize}

\subsection*{رفلکس‌ها}
\begin{itemize}
  \item مورو: از بدو تولد، تا ۳–۶ ماهگی.
  \item گرسپ: از بدو تولد، تا ۳–۴ ماهگی.
  \item \en{Rooting}: از بدو تولد، تا ۳–۴ ماهگی.
  \item گالانت: از بدو تولد، تا ۴ ماهگی.
  \item \en{Placing}: از بدو تولد، تا ۲–۶ ماهگی.
\end{itemize}
شمارهٔ دست‌نویس پایین صفحه: ۱.

\clearpage
\pagehead{۲}
\subsection*{ادامهٔ رفلکس‌ها}
\begin{itemize}
  \item آسمتریک تانیک نک: از بدو تولد، تا ۳ ماهگی.
  \item پاراشوت: \marktext{از ۸–۱۰ ماهگی، تا پایان عمر}.
  \item بابینسکی: از بدو تولد، تا \marktext{۱۲–۱۸ ماهگی}.
\end{itemize}

\subsection*{تون و حرکات نوزاد}
\begin{itemize}
  \item جنین پس از ۲۸ هفتگی: \en{lie}؛ اکستنشن اندام‌ها (در حالت استراحت).
  \item سپس در ۳۲ هفتگی: تون و فلکشن اندام تحتانی افزایش می‌یابد.
  \item در ۳۶ هفتگی: اندام تحتانی فلکس، فوقانی اکستنت.
  \item جنین ترم: هر دو اندام‌ها فلکس.
  \item \en{Recoil}: وقتی اندام‌ها را باز و اکستنت می‌کنیم، می‌تواند فلکس کند؛ در جنین‌های پره‌مچور وجود ندارد.
  \item حرکات نوزاد پره‌مچور کند است.
\end{itemize}

\subsection*{معاینهٔ نورولوژیک}
معاینهٔ نورولوژیک نوزاد:
\begin{latin}
\begin{itemize}
  \item Mental Status
  \item Reflexes
  \item Motor examination and tone
\end{itemize}
\end{latin}
معاینهٔ نورولوژیک کودک:
\begin{latin}
\begin{itemize}
  \item Mental status and developmental exam
  \item Cranial nerve examination
  \item Motor exam, Sensory exam, DTR
\end{itemize}
\end{latin}
\begin{itemize}
  \item \en{Alertness}: فعالیت خودبه‌خودی، تغذیه، فیکس اند فالو، والدین، محرک‌ها.
  \item عملکرد زبانی: ساده‌ترین راه ارزیابی هوش.
  \begin{itemize}
    \item بخش شناختی: فهم جمله.
    \item بخش بیانی: تکلم جمله.
  \end{itemize}
  \item اختلالات زبان ناشی از اختلالات نیمکره‌های مغزی: \marktext{آفازی}.
\end{itemize}
شمارهٔ دست‌نویس پایین صفحه: ۲.

\clearpage
\pagehead{۳}
\subsection*{آفازی}
\begin{itemize}
  \item آفازی بروکا: قدامی، اکسپرسیو.
  \begin{latin}
  $\rightarrow$ Spare, non-fluent language.
  \end{latin}
  \item آفازی ورنیکه: خلفی، رسپتیو.
  \begin{latin}
  $\rightarrow$ Inability to understand, \marktext{speech is fluent but nonsensical}.
  \end{latin}
  \item آفازی گلوبال: هر دو.
\end{itemize}

\subsection*{ارزیابی اعصاب کرانیال}
\begin{itemize}
  \item عصب \en{II}: نوزاد ترم، نور و صورت انسان را دنبال می‌کند.
  \begin{itemize}
    \item پاسخ کاهش‌یافتهٔ مردمک به آسیب: مسیر بینایی قدامی (رتین، عصب، کیاسم).
    \item \en{Unilateral lesions $\rightarrow$ Swinging Flashlight test}.
    \item \en{\marktext{Marcus Gunn pupil} $\rightarrow$ afferent pupillary defect}؛ دیلاتاسیون غیرقرینه.
    \item \en{\marktext{Horner Syndrome} (ptosis, miosis, unilateral facial anhidrosis)}: نقص عصب‌دهی سمپاتیک به چشم یک طرف.
    \item آنیزوکوریا (مردمک‌های غیرقرینه): در \uncertain{نارسی مغز} دیده می‌شود.
    \item آسیب‌های مسیر خلفی بینایی: رفلکس‌های مردمک \en{NL} است؛ نابینایی دارند.
  \end{itemize}
  \item عصب‌های \en{III, IV, VI}: حرکات چشم.
  \begin{itemize}
    \item شیرخواران: \en{Oculocephalic reflex (\marktext{Doll's eye maneuver})}.
    \item اگر ساقهٔ مغز سالم باشد، چرخاندن جمجمه به یک سمت، منجر به چرخیدن چشم به سمت مقابل می‌شود.
    \item اگر خراب باشد، چشم‌ها همان‌جا می‌ماند.
    \item فلج یک‌طرفهٔ عصب \en{III} (اکولوموتور): چشم به پایین و خارج می‌رود؛ پتوز + مردمک میدریاتیک.
    \item فلج عصب \en{IV} (تروکلئار): ضعف حرکت رو به پایین + دوبینی عمودی.
  \end{itemize}
\end{itemize}
شمارهٔ دست‌نویس پایین صفحه: ۳.

\clearpage
\pagehead{۴}
\subsection*{ادامهٔ ارزیابی اعصاب کرانیال}
\begin{itemize}
  \item عصب \en{VI} (ابدوسنس) اگر فلج شود:
  \begin{itemize}
    \item ناتوانی در حرکت چشم به خارج.
    \item چون عصب طولانی است، افزایش \en{ICP} اول آن را درگیر می‌کند.
  \end{itemize}
  \item عصب \en{V} (تری‌ژمینال):
  \begin{itemize}
    \item \en{\marktext{Corneal reflex}}:
    \begin{latin}
    \en{V} $\rightarrow$ Ophthalmic division (sensory arch)\\
    \en{VII} (motor arch)
    \end{latin}
    \item لمس چانه: رفلکس روتینگ.
  \end{itemize}
  \item عصب \en{VII}: مشاهدهٔ صورت حین استراحت؛ \en{Bell palsy}.
  \begin{itemize}
    \item ضعف فقط در پایین صورت (دهان):
    \begin{latin}
    Contralateral lesion of the upper motor neuron.
    \end{latin}
  \end{itemize}
  \item عصب \en{VIII}: ناشنوایی، تینیتوس، سرگیجه.
  \begin{itemize}
    \item جزء وستیبولار: سرگیجه، تهوع، استفراغ، نیستاگموس.
  \end{itemize}
  \item عصب‌های \en{IX, X}:
  \begin{itemize}
    \item رفلکس \en{gag} در هر سن وجود دارد.
    \item نبودش: آسیب به ساقهٔ مغز، ۹ و ۱۰.
    \item آسیب \en{X}: دیسفاژی، ساکینگ ضعیف.
  \end{itemize}
  \item عصب \en{XI}: عملکرد عضلهٔ تراپزیوس و استرنوکلیدوماستوئید.
  \begin{itemize}
    \item آسیب به \en{XI}: \en{Head tilt, dropping of the shoulder}.
  \end{itemize}
  \item عصب \en{XII}:
  \begin{itemize}
    \item آتروفی و فاسیکولاسیون لترال زبان.
    \item در خواب بهتر دیده می‌شود.
    \item زبان به طرف عصب ضعیف منحرف می‌شود.
  \end{itemize}
\end{itemize}
شمارهٔ دست‌نویس پایین صفحه: ۴.

\clearpage
\pagehead{۵}
\subsection*{تون، بالک و قدرت عضله}
\begin{itemize}
  \item تون: مقاومت دینامیک عضله به حرکات \en{passive}.
  \begin{itemize}
    \item ضایعات \en{LMN} و ضایعات مخچه: \marktext{هایپوتونی}.
    \item ضایعات \en{UMN} و اکستراپیرامیدال (بازال گانگلیا): \marktext{هایپرتونی}.
    \item \en{UMN}: اسپاستیسته؛ اکستراپیرامیدال: ریجیدیتی.
  \end{itemize}
  \item بالک عضله:
  \begin{itemize}
    \item ضایعات \en{LMN}، نوروماسکولار: کاهش بالک.
    \item \marktext{مایوتونیا کانجنیتا}: افزایش بالک.
    \item دیستروفی دوشن: \marktext{سودوهایپرتروفی اندام تحتانی}.
  \end{itemize}
  \item \en{Gower Sign}: برای بلند شدن از خودش بالا می‌رود؛ ضعف شدید عضلات \marktext{پروگزیمال}.
  \item فاسیکولاسیون عضلات: دنرواسیون ناشی از آسیب شاخ قدامی یا اعصاب محیطی.
  \item آتاکسی (نبود کوردینیشن): آسیب به مسیرهای مخچه‌ای؛ دیسمتری (نقص در تخمین فاصله).
  \item \marktext{در ۴ سالگی کودک باید بتواند روی انگشت و پاشنه راه برود.}
  \item آسیب‌های مخچه اختلال گیت می‌دهد.
  \item آسیب مسیر کورتیکواسپاینال:
  \begin{latin}
  Stiff, \marktext{scissoring} gate and toe walking.
  \end{latin}
  حاشیهٔ دست‌نویس دربارهٔ \en{scissoring}: حین راه رفتن از خط وسط رد می‌شود.
\end{itemize}

\subsection*{نمره‌دهی \en{DTR}}
\begin{itemize}
  \item \en{0}: نبود رفلکس.
  \item \en{1+}: \en{Trace}.
  \item \en{2+}: \en{NL}.
  \item \en{3+}: افزایش‌یافته.
  \item \en{4+}: با کلونوس.
\end{itemize}
شمارهٔ دست‌نویس پایین صفحه: ۵.

\clearpage
\pagehead{۶}
\subsection*{رفلکس‌ها و حس}
\begin{itemize}
  \item رفلکس‌های \en{DTR} در بیماری‌های \en{UMN} افزایش و با \en{LMN} کاهش می‌یابند.
  \item رفلکس بابینسکی: حرکت رو به بالای انگشت بزرگ در تحریک لترال پا.
  \begin{itemize}
    \item نشان آسیب به مسیر \marktext{کورتیکواسپاینال}.
    \item در نوزادان کاربرد ندارد، مگر اینکه غیرقرینه باشد.
    \item در \marktext{۱۲–۱۸ ماهگی} قابل انجام؛ \en{NL}: انگشت فلکس می‌شود.
  \end{itemize}
  \item تحریک یک اندام در شیرخوار: باید به آن نگاه کند یا گریه کند.
  \item \en{Stereognosis}: شناسایی جسم قرار داده‌شده در دست.
  \item \en{Graphesthesia}: شناسایی عدد نوشته‌شده روی دست.
  \item \en{Two point discrimination}: توانایی افتراق دو جسم تیز.
\end{itemize}

\subsection*{\en{CSF}}
\begin{itemize}
  \item افزایش \en{PMN} و کاهش گلوکز: عفونت باکتریال، آبسهٔ مغزی، عفونت انگلی.
  \item افزایش لنفوسیت و کاهش گلوکز: \en{TB}، سارکوئیدوز، عفونت قارچی.
  \item افزایش لنفوسیت و گلوکز \en{NL}: مننژیت وایرال، \en{ADEM}، \en{MS}، واسکولیت، تومور.
  \item افزایش پروتئین: عفونت، گیلن‌باره، …
  \item گلوکز \en{NL} در \en{LP}: در نوزاد \en{$>40$}، شیرخوار \en{$>60$}؛ بالای \en{$\frac{2}{3}$} گلوکز خون.
  \item پروتئین \en{NL} در \en{CSF}: نوزاد ۶۵–۱۵۰، شیرخوار \en{$<40$}.
\end{itemize}

\subsection*{\en{EEG}}
\begin{itemize}
  \item \en{persistent focal slowing}: \en{Tumor}، استروک، آبسه.
  \item \en{Diffuse slow wave activity}: متابولیک انسفالوپاتی، افزایش \en{ICP}.
\end{itemize}
شمارهٔ دست‌نویس پایین صفحه: ۶.
\ifPedNeuroEvalStandalone
\end{document}
\fi
